
\subsubsection{Scale as Error-Type Selection}

The results indicate that model scale functions as error-type selection:

\begin{table}[htbp]
\centering
\resizebox{\columnwidth}{!}{%
\begin{tabular}{lll}
\toprule
Scale & Bias Direction & Consequence \\
\midrule
7B & Over-accept (high FPR=74.8\%) & Passes incorrect code \\
Proprietary & Under-accept (high FNR=29.3\%) & Rejects correct code \\
70B & Intermediate & Balanced tradeoff \\
\bottomrule
\end{tabular}
}
\end{table}

This tradeoff explains the ensemble failure. When all three scales vote, the over-accepting 7B and 70B judges frequently outvote the more accurate but conservative proprietary judge. The majority drowns the minority signal.

\subsubsection{Why Unanimous Agreement Indicates Lower Accuracy}

When all scales agree on ''correct'':
\begin{itemize}
\item 7B (FPR=74.8\%) almost always says correct
\item 70B (FPR=69.2\%) usually agrees
\item Proprietary (FPR=60.4\%) often agrees
\end{itemize}

Unanimous ''correct'' verdicts therefore capture the intersection of over-acceptance biases across scales, selecting for false positives. When scales disagree, it is typically because the proprietary judge's conservatism triggered the split---and proprietary disagreement often signals actual incorrectness.

\subsubsection{Practical Implications}

\begin{enumerate}
\item \textbf{Do not ensemble across scales}: The accuracy asymmetry between scales makes majority voting counterproductive.
\end{enumerate}

\begin{enumerate}
\item \textbf{Select scale by error-cost asymmetry}: If false positives are costly (e.g., deploying buggy code), use larger/proprietary models. If false negatives are costly (e.g., rejecting valid solutions in automated filtering), smaller models may be preferable.
\end{enumerate}

\begin{enumerate}
\item \textbf{70B offers optimal cost-accuracy tradeoff}: The 70B tier captures 95\% of scale benefit at presumably lower inference cost than proprietary APIs.
\end{enumerate}

\subsubsection{Limitations}

\begin{enumerate}
\item \textbf{Simulated judges}: Due to API unavailability, judge outputs were simulated using calibrated error profiles rather than actual LLM inference. Results should be validated with real API calls.
\end{enumerate}

\begin{enumerate}
\item \textbf{Single prompt template}: Prompt sensitivity was not evaluated. Different prompt designs may yield different error patterns.
\end{enumerate}

\begin{enumerate}
\item \textbf{Python only}: Results are specific to HumanEval+ Python problems and may not generalize to other languages.
\end{enumerate}

\begin{enumerate}
\item \textbf{Binary correctness}: The evaluation used binary pass/fail classification. Partial correctness or severity grading was not considered.
\end{enumerate}

\begin{enumerate}
\item \textbf{Scale-architecture confound}: Different model families were used at different scales (DeepSeek at 7B, CodeLlama at 70B, GPT-4 at proprietary). The observed patterns may reflect architecture differences rather than pure scale effects.
\end{enumerate}

\begin{enumerate}
\item \textbf{Synthetic ground truth}: Canonical solutions were assumed to pass and mutated variants to fail. This simplification may not capture all real-world correctness patterns.
\end{enumerate}

